\chapter{Experimental setup}\label{ch:setup}

\section{Data, baselines and training}\label{sec:training}

We use DIV2K \cite{agustsson2017div2k}, with 800 training images and 100
validation images. Every reported result comes from the validation images.
Training reads from a cache of 32 fixed $128\times128$ crops per training
image, 25,600 tiles in all, with random flips and rotations. The degradation
runs on the fly, so each epoch sees new corruptions. For reproducibility we
also exported a fixed set of 5,000 pairs generated with seed 1234.

The baselines are an RGB DnCNN (17 layers, width 64, 558,403 parameters) and
FFDNet in the author's colour configuration (852,108 parameters), in two
variants. \emph{Oracle} FFDNet gets the true sampled noise level. No real
user would have that, so it serves as an upper reference. \emph{Blind}
FFDNet gets a constant map of $27.5/255$ and sees exactly what the other
models see. We train all four from scratch with the same settings
(Table~\ref{tab:hparams}). Each model trained once with an unfixed seed, on a
Tesla T4 in Kaggle sessions that resume from checkpoints.

\begin{table}[htbp]
	\centering
	\caption{Training configuration, the same for every method.}
	\label{tab:hparams}
	\small
	\begin{tabular}{@{}ll@{}}
		\toprule
		Setting & Value \\
		\midrule
		Optimiser and schedule & AdamW \cite{loshchilov2019adamw}, $2\times10^{-4}$ cosine \cite{loshchilov2017sgdr} to $10^{-6}$ \\
		Updates and batch      & 200,000; 16 crops of $128\times128$ \\
		Loss                   & $\mathcal{L}_\text{char}+0.05\,\mathcal{L}_\text{freq}$ (+ entropy for \fcsg{}) \\
		Precision              & FP16 mixed precision \cite{micikevicius2018amp} \\
		Hardware               & NVIDIA Tesla T4 (Kaggle) \\
		\bottomrule
	\end{tabular}
\end{table}

\section{Evaluation protocol}\label{sec:evalproto}

We degrade each of the 100 validation images once, with its index as the
seed, so all models see identical inputs. Inference runs on $256\times256$
tiles with 32 pixels of overlap, averaged where tiles meet. During training
we monitored PSNR on 16 fixed crops. That number does not compare with
full-image results. For \fcsg{} it is 30.59\dB{} on crops and 27.24\dB{} on
full images.

\section{Verification and baseline reproduction}\label{sec:verification}

Before training, a notebook on the same T4 ran the checks in
Table~\ref{tab:gates} and stopped at the first failure. The overfit check
earns its place. If a model cannot memorise one image in 200 steps it has a
bug, and finding that out in seconds beats finding it after ten hours of
training.

\begin{table}[htbp]
	\centering
	\caption{Checks run before training.}
	\label{tab:gates}
	\small
	\begin{tabular}{@{}lll@{}}
		\toprule
		Check & Criterion & Measured \\
		\midrule
		FFDNet weights, CBSD68 $\sigma{=}25$ & within 0.5\dB{} of 31.21 & 31.220\dB \\
		Band reconstruction & $x_L{+}x_M{+}x_H = x$ & max error $6.0\times10^{-7}$ \\
		Parameters / GFLOPs at $256^2$ & $<5$M / $<10$ & 97,088 / 9.39 \\
		Memory, forward + backward & $< 4$ GB & 0.60 GB \\
		One-image overfit, 200 steps & MSE $<0.01$, $<\frac14$ initial & ratio 0.18 \\
		Checkpoint restore, reproducible data & identical & passed \\
		\bottomrule
	\end{tabular}
\end{table}

To check our FFDNet code against a published number, we ran the author's
pretrained colour weights \cite{kair} on the 68 CBSD68 images
\cite{martin2001bsd} with the author's protocol. We got 31.220\dB{} against
31.21\dB{} in the paper, a gap of 0.01\dB{}. To test the training loop end to
end, we trained DnCNN on clipped Gaussian noise at $\sigma=25$. It reached
32.65\dB{} on DIV2K validation and 31.04\dB{} on CBSD68.
